\chapter{Provenance des figures et contact}
\label{ann:figure-provenance}

\section{Référence scientifique et niveaux de vérification}
Sources documentaires du 9 octobre 2026, PDF reconstruit et revu le 10 octobre
2026. La \textbf{provenance scientifique}
est rattachée au commit de référence
\href{https://github.com/halhoulmhamed-droid/pde-image-denoising/commit/ae03109cdb2150b760e5a4b8bcffb56bba1b98e0}{\texttt{ae03109cdb21}},
qui scelle les fichiers de code et de résultats consultés :
\begin{center}
\nolinkurl{ae03109cdb2150b760e5a4b8bcffb56bba1b98e0}
\end{center}
Ce SHA n'est pas celui de la présente édition PDF. Il identifie l'archive
consultée, non une preuve autonome du commit exécuté lors de l'expérience :
le manifeste historique ne consigne pas cet identifiant d'exécution.
Le lien du commit documentaire ne peut être fixé qu'après création de l'édition.

Les contrôles de cette édition sont administratifs : concordance des chemins,
tailles, empreintes et références archivées ; lecture de la configuration et
du code de tracé ; comparaison des empreintes de tableaux avec les empreintes
déjà enregistrées. Aucun solveur, calcul de métrique ou agrégation statistique
n'a été relancé. Les rapports de vérification historiques restent des preuves
de leurs exécutions antérieures, distinctes de ces contrôles documentaires.
L'accord des fichiers ne prouve pas à lui seul leur exécution scientifique.

Le registre détaillé
\href{https://github.com/halhoulmhamed-droid/pde-image-denoising/blob/main/docs/figure_provenance.md}{\texttt{docs/figure\_provenance.md}}
donne les SHA-256, les fonctions génératrices, les usages et les limites
par figure. Les liens de données et de code ci-dessous sont des permaliens
au commit de provenance scientifique ; le registre et les PDF de cette édition
utilisent au contraire les chemins courants de la branche \texttt{main}.

\section{Correspondance des illustrations}
\noindent
\begin{longtable}{@{}p{0.9cm}p{6.1cm}p{7.7cm}@{}}
\toprule
ID & Source graphique archivée & Usage et type de provenance \\
\midrule
\endfirsthead
\toprule
ID & Source graphique archivée & Usage et type de provenance \\
\midrule
\endhead
F01 & \href{https://github.com/halhoulmhamed-droid/pde-image-denoising/blob/ae03109cdb2150b760e5a4b8bcffb56bba1b98e0/results/extended/figures/comparison.pdf}{Comparaison des cinq images (PDF)} & Figure~\ref{fig:comparison} ; slides : comparaison. Résultat individuel, source A. \\
F02 & \href{https://github.com/halhoulmhamed-droid/pde-image-denoising/blob/ae03109cdb2150b760e5a4b8bcffb56bba1b98e0/results/extended/figures/intensity_profile.pdf}{Profil horizontal (PDF)} & Figure~\ref{fig:profile} ; slides : profil. Résultat individuel, ligne 42, source A. \\
F03 & \href{https://github.com/halhoulmhamed-droid/pde-image-denoising/blob/ae03109cdb2150b760e5a4b8bcffb56bba1b98e0/results/extended/figures/gradient_maps.pdf}{Cartes de gradient (PDF)} & Figure~\ref{fig:gradients} ; slides : gradients. Transformation d'affichage du résultat individuel, source A ; pas une métrique de contours. \\
F04 & \href{https://github.com/halhoulmhamed-droid/pde-image-denoising/blob/ae03109cdb2150b760e5a4b8bcffb56bba1b98e0/results/extended/figures/metric_curves_geometric_shapes_sigma0p025.pdf}{Courbes : formes, $\sigma=0{,}025$ (PDF)} & Figure~\ref{fig:curves025}, panneau supérieur. Agrégat, source B. \\
F05 & \href{https://github.com/halhoulmhamed-droid/pde-image-denoising/blob/ae03109cdb2150b760e5a4b8bcffb56bba1b98e0/results/extended/figures/metric_curves_ramps_and_edges_sigma0p025.pdf}{Courbes : rampes, $\sigma=0{,}025$ (PDF)} & Figure~\ref{fig:curves025}, panneau inférieur. Agrégat, source B. \\
F06 & \href{https://github.com/halhoulmhamed-droid/pde-image-denoising/blob/ae03109cdb2150b760e5a4b8bcffb56bba1b98e0/results/extended/figures/metric_curves_geometric_shapes_sigma0p05.pdf}{Courbes : formes, $\sigma=0{,}05$ (PDF)} & Figure~\ref{fig:curves05}, panneau supérieur. Agrégat, source B. \\
F07 & \href{https://github.com/halhoulmhamed-droid/pde-image-denoising/blob/ae03109cdb2150b760e5a4b8bcffb56bba1b98e0/results/extended/figures/metric_curves_ramps_and_edges_sigma0p05.pdf}{Courbes : rampes, $\sigma=0{,}05$ (PDF)} & Figure~\ref{fig:curves05}, panneau inférieur. Agrégat, source B. \\
F08 & \href{https://github.com/halhoulmhamed-droid/pde-image-denoising/blob/ae03109cdb2150b760e5a4b8bcffb56bba1b98e0/results/extended/figures/metric_curves_geometric_shapes_sigma0p1.pdf}{Courbes : formes, $\sigma=0{,}10$ (PDF)} & Figure~\ref{fig:curves10}, panneau supérieur ; slides : PSNR puis SSIM, recadrés. Agrégat, source B. \\
F09 & \href{https://github.com/halhoulmhamed-droid/pde-image-denoising/blob/ae03109cdb2150b760e5a4b8bcffb56bba1b98e0/results/extended/figures/metric_curves_ramps_and_edges_sigma0p1.pdf}{Courbes : rampes, $\sigma=0{,}10$ (PDF)} & Figure~\ref{fig:curves10}, panneau inférieur ; slides : PSNR puis SSIM, recadrés. Agrégat, source B. \\
F10 & \href{https://github.com/halhoulmhamed-droid/pde-image-denoising/blob/ae03109cdb2150b760e5a4b8bcffb56bba1b98e0/results/extended/figures/k_sensitivity_geometric_shapes_sigma0p1.pdf}{Sensibilité à $K$ : formes, $\sigma=0{,}10$ (PDF)} & Figure~\ref{fig:k-sensitivity}, panneau supérieur ; slides : PSNR puis SSIM, recadrés. Agrégat, source B. \\
F11 & \href{https://github.com/halhoulmhamed-droid/pde-image-denoising/blob/ae03109cdb2150b760e5a4b8bcffb56bba1b98e0/results/extended/figures/k_sensitivity_ramps_and_edges_sigma0p1.pdf}{Sensibilité à $K$ : rampes, $\sigma=0{,}10$ (PDF)} & Figure~\ref{fig:k-sensitivity}, panneau inférieur ; réserves : PSNR puis SSIM, recadrés. Agrégat, source B. \\
\bottomrule
\end{longtable}

\subsection*{Source A : cas individuel archivé}
Les cinq états sont l'image propre, l'observation bruitée, la chaleur et les
deux PM avec $K=0{,}10$, pour \texttt{geometric\_shapes}, $\sigma=0{,}10$,
graine 0, $n=20$, $t=4$, $dx=dy=1$ et $dt=0{,}20$.
Le \href{https://github.com/halhoulmhamed-droid/pde-image-denoising/blob/ae03109cdb2150b760e5a4b8bcffb56bba1b98e0/results/extended/arrays/representative_outputs.npz}{NPZ des états représentatifs}
archive l'objet en mémoire passé historiquement aux fonctions
\href{https://github.com/halhoulmhamed-droid/pde-image-denoising/blob/ae03109cdb2150b760e5a4b8bcffb56bba1b98e0/src/pde_image_denoising/plotting.py}{\texttt{comparison\_figure}, \texttt{profile\_figure} et \texttt{gradient\_figure}}.
La chaîne lue dans le code sauvegarde cet objet avant le tracé ; les fonctions
ne relisent pas le NPZ pour créer ces trois images. Les empreintes des cinq
tableaux sont comparées aux champs déjà enregistrés du
\href{https://github.com/halhoulmhamed-droid/pde-image-denoising/blob/ae03109cdb2150b760e5a4b8bcffb56bba1b98e0/results/extended/metrics.csv}{CSV de métriques},
sans calculer une nouvelle métrique. Les cartes de gradient appliquent la
transformation d'affichage documentée, sans mesure nouvelle de contours.

\subsection*{Source B : agrégats déjà calculés}
Les courbes sont issues de
\href{https://github.com/halhoulmhamed-droid/pde-image-denoising/blob/ae03109cdb2150b760e5a4b8bcffb56bba1b98e0/results/extended/summary.csv}{\texttt{summary.csv}},
lu historiquement par
\href{https://github.com/halhoulmhamed-droid/pde-image-denoising/blob/ae03109cdb2150b760e5a4b8bcffb56bba1b98e0/src/pde_image_denoising/extended_analysis.py}{\texttt{metric\_curves\_figure} et \texttt{k\_sensitivity\_figure}}.
Elles montrent les moyennes et écarts-types échantillonnaux sur les graines
0 à 9 ; bandes et barres sont un écart-type, pas un intervalle de confiance.
Les courbes utilisent les huit checkpoints archivés ; la sensibilité montre
$n=5,20,80$. Les deux conductances et les quatre $K$ sont conservés.
Les recadrages des slides isolent les panneaux PSNR et SSIM ; leurs légendes
latérales sont une recomposition de présentation, non un nouveau résultat.

\subsection*{Configuration, intégrité et limites}
Les liens permanents vers
\href{https://github.com/halhoulmhamed-droid/pde-image-denoising/blob/ae03109cdb2150b760e5a4b8bcffb56bba1b98e0/results/extended/config_used.json}{la configuration utilisée},
\href{https://github.com/halhoulmhamed-droid/pde-image-denoising/blob/ae03109cdb2150b760e5a4b8bcffb56bba1b98e0/results/extended/manifest.json}{le manifeste SHA-256},
\href{https://github.com/halhoulmhamed-droid/pde-image-denoising/blob/ae03109cdb2150b760e5a4b8bcffb56bba1b98e0/results/extended/verification.json}{la vérification stricte historique}
et \href{https://github.com/halhoulmhamed-droid/pde-image-denoising/blob/ae03109cdb2150b760e5a4b8bcffb56bba1b98e0/src/pde_image_denoising/experiments.py}{le code de la chaîne}
complètent cette correspondance. Les logos relèvent séparément des
\href{https://github.com/halhoulmhamed-droid/pde-image-denoising/blob/ae03109cdb2150b760e5a4b8bcffb56bba1b98e0/assets/logos/SOURCES.md}{sources institutionnelles archivées}.
Les marqueurs TikZ des légendes sont pédagogiques ; ils n'ajoutent aucune
donnée expérimentale. La lecture et les empreintes ne constituent ni une
reproduction relancée, ni une authentification de l'exécution scientifique,
ni une preuve de l'origine des pixels raster. La revue du rendu de cette
édition reste distincte de la provenance matérielle des archives.

\clearpage
\section{Contact}
\label{ann:contact-halhoul}
\noindent\textbf{Halhoul Mhamed}\par
\smallskip
\noindent
\href{mailto:halhoul.mhamed@etu.uae.ac.ma}{\nolinkurl{halhoul.mhamed@etu.uae.ac.ma}}\par
\noindent
\href{mailto:halhoulmhamed@gmail.com}{\nolinkurl{halhoulmhamed@gmail.com}}\par
\medskip
\noindent
\href{https://github.com/halhoulmhamed-droid/pde-image-denoising}{GitHub : \nolinkurl{halhoulmhamed-droid/pde-image-denoising}}
